\documentclass[11pt,a4paper]{article}
\usepackage{forjatex-manual}

\hypersetup{
  colorlinks=true,
  linkcolor=fjblue,
  urlcolor=fjblue,
  pdftitle={Manual do Usuário - ForjaBib},
  pdfauthor={Jean Pimenta / Arkhé Code}
}

\begin{document}

\ForjaHeader{../../forjabib/src-tauri/icons/128x128@2x.png}
  {Ecossistema ForjaTeX App Suite $\cdot$ GNU GPL v3.0}
  {Manual do Usuário --- ForjaBib}
  {Gerenciador Visual de Bibliografias BibTeX e BibLaTeX (ABNT)}
  {Versão 1.0.5}
  {Jean Pimenta $\cdot$ \href{https://forjatex.com.br}{\color{fjsky}forjatex.com.br} $\cdot$ \href{https://github.com/jeanjpimenta/forjatex-suite}{\color{fjsky}GitHub}}

\ForjaTOC{Sumário}

\begin{fjcallout}[Arquivo .bib como Fonte Única da Verdade]
O \textbf{ForjaBib} edita diretamente arquivos \code{.bib} limpos e compatíveis com Git, sem banco proprietário oculto e preservando macros \LaTeX{} como \code{\textbackslash apj} e fórmulas matemáticas \code{\$...\$}.
\end{fjcallout}

\section{Introdução}
O \textbf{ForjaBib} é o gerenciador de referências bibliográficas do ecossistema \textbf{ForjaTeX App Suite}, projetado para eliminar os erros clássicos de compilação no \code{biber} e \code{bibtex}, com foco especial no padrão brasileiro \textbf{ABNT (\code{biblatex-abnt})}.

\section{Importação Automática (DOI, arXiv e ISBN)}
\begin{enumerate}
  \item \textbf{DOI (CrossRef):} cole um DOI (ex.: \code{10.1103/PhysRev.76.769}) e clique em \textbf{Buscar}.
  \item \textbf{arXiv ID:} informe um identificador arXiv (ex.: \code{1412.6980}); o ForjaBib cria automaticamente uma entrada \code{@misc} com \code{eprint}, \code{archivePrefix = \{arXiv\}} e \code{primaryClass}.
  \item \textbf{ISBN (OpenLibrary):} informe o ISBN de 10 ou 13 dígitos para preencher livros automaticamente.
\end{enumerate}

\section{Linter Pre-Flight e Proteção ABNT}
Antes mesmo de compilar no \textbf{ForjaTeX IDE}, o ForjaBib verifica:
\begin{itemize}
  \item \textbf{Chaves duplicadas:} detecta conflitos de \code{@chave} e gera chaves únicas padronizadas (\code{autor2026palavra}).
  \item \textbf{Caracteres não escapados e chaves desbalanceadas:} alerta sobre \code{\&}, \code{\%}, \code{\#} sem barra invertida ou chaves \code{\{\}} abertas.
  \item \textbf{Conflito DOI + URL:} identifica registros que repetem o link do DOI em \code{url} e remove em 1 clique.
  \item \textbf{Proteger Siglas (\code{\{ABC\}}):} envolve siglas como \code{\{COVID-19\}} e \code{\{USA\}} sem colocar chaves duplas no título inteiro.
  \item \textbf{Autor Corporativo (\code{\{\{...\}\}}):} envolve instituições como \code{\{\{IBGE\}\}} para que o \code{biblatex-abnt} não inverta o nome.
  \item \textbf{Validação \code{urldate}:} avisa quando uma referência com \code{url} está sem data de acesso exigida pela NBR 6023.
\end{itemize}

\section{Auditoria .tex e Prévia Formatada}
\begin{itemize}
  \item \textbf{Prévia instantânea:} visualize a referência formatada em \textbf{ABNT}, \textbf{IEEE}, \textbf{APA} ou \textbf{Author-Year}.
  \item \textbf{Auditar .tex:} selecione um arquivo \code{.tex} do seu projeto para descobrir citações faltantes no \code{.bib} (que gerariam \code{[?]}) e entradas órfãs nunca citadas.
  \item \textbf{Detector de Duplicatas:} compara por DOI e título semelhante para mesclar/remover cópias.
\end{itemize}

\section{Histórico de Versões (Changelog)}
\begin{itemize}
  \item \textbf{v1.0.5:} Novo 3º motor \textbf{CSL (\code{citation-style-language})} com importador local de arquivos \code{.csl}, download direto do repositório oficial Zotero/CSL, prévia exata via \code{citeproc-js}, teste em PDF e extração de pistas \code{.csl} para o Criador BibLaTeX.
  \item \textbf{v1.0.4:} Novo modo \textbf{Criador de Estilo} por aba para personalizar estilos \code{biblatex} e \code{natbib} visualmente (\code{giveninits}, \code{maxbibnames}, \code{\textbackslash DeclareNameAlias}, \code{\textbackslash DeclareFieldFormat}, \code{\textbackslash newunitpunct}), com prévia ao vivo, exportação de \code{biblatex.cfg} e compilação PDF de teste.
  \item \textbf{v1.0.3:} Lixeira de referências com confirmação, retenção automática de 30 dias e restauração para o acervo, além de gravação otimizada em segundo plano.
  \item \textbf{v1.0.2:} Suporte dedicado ao motor \code{natbib} (BibTeX clássico) por aba com conversão automática de tipos e campos (\code{journaltitle}, \code{location}, \code{date}), atalhos \code{\textbackslash citep}/\code{\textbackslash citet}, prévia no texto e gerador de preâmbulo \code{natbib}/\code{abntex2cite}.
  \item \textbf{v1.0.1:} Abas independentes e renomeáveis com biblioteca global compartilhada por checkboxes à esquerda, suporte completo aos 35 tipos (\code{@type}) do BibTeX e BibLaTeX e 18 estilos de prévia formatada com busca em lista suspensa.
  \item \textbf{v1.0.0:} Lançamento inicial com edição bidirecional \code{.bib}, importador DOI/arXiv/ISBN, linter pre-flight, proteção de siglas e autor corporativo, prévia ABNT/IEEE/APA, detector de duplicatas e auditoria contra arquivos \code{.tex}.
\end{itemize}

\end{document}

